\chapter{Methodology and Implementation}
\label{ch:methodology}

\section{Chapter Overview}
\label{sec:methodology_overview}

This chapter describes the methodology used to design, implement, and
evaluate the proposed broadcast football video analysis system. Building on
the background and related work discussed in Chapter~\ref{ch:background},
the aim is to construct an open and reproducible pipeline that takes a single
broadcast football video as input and produces both visual and quantitative
tactical outputs.

The chapter begins by defining the requirements and design objectives that
guided the system. It then presents the overall system architecture, including
the main processing stages, intermediate representations, and final outputs.
The following sections describe the datasets and training procedures used for
the detection and pitch keypoint models, before explaining the main pipeline
components: object and ball detection, multi-object tracking, team
identification, pitch registration, and tactical output generation.

Particular attention is given to the design choices made at each stage, since
the performance of later components depends directly on the reliability of
earlier ones. For example, ball detection affects possession and pass
estimation, while pitch registration affects heatmaps, speed calculation, and
offside candidate detection. The chapter therefore describes not only what was
implemented, but also why each method was selected and how errors are managed
where possible.

The final sections describe the evaluation methodology and reproducibility
measures. These include component-level evaluation, end-to-end pipeline
evaluation, ablation studies, runtime measurement, and error analysis. Results
from these evaluations are reported in Chapters~\ref{ch:results}
and~\ref{ch:evaluation}.

\section{Requirements and Design Objectives}
\label{sec:requirements_design_objectives}

The methodology of the system was guided by the overall aim of producing an
open and reproducible pipeline for tactical football analysis from broadcast
video alone. This required the system to recover player, ball, team, and pitch
information from a single moving-camera feed, and to transform these
intermediate outputs into analytically useful representations such as heatmaps,
possession statistics, pass events, and candidate offside situations.

To ensure that the design remained aligned with this aim, the system was
specified in terms of functional requirements, non-functional requirements, and
success criteria. These requirements were used both to guide implementation
choices and to define the basis for later evaluation.

\subsection{Functional Requirements}
\label{subsec:functional_requirements}

The functional requirements describe the core capabilities that the system was
expected to provide.

\begin{itemize}
    \item \textbf{FR1: Object detection.} The system must detect the main
    entities required for football analysis, namely players, goalkeepers,
    referees, and the ball, from broadcast frames.

    \item \textbf{FR2: Temporal tracking.} The system must maintain player
    identities across frames so that movement-based analytics can be computed
    over time.

    \item \textbf{FR3: Team identification.} The system must assign players to
    teams without requiring match-specific labelled training data for each new
    fixture.

    \item \textbf{FR4: Pitch registration.} The system must estimate a mapping
    from image coordinates to a canonical pitch model so that image-space
    detections can be analysed in real-world pitch space.

    \item \textbf{FR5: Tactical output generation.} The system must derive
    higher-level outputs from the tracked and registered data, including player
    heatmaps, speed and distance estimates, team possession statistics, pass
    events, and candidate offside situations.

    \item \textbf{FR6: Visual output generation.} The system must produce
    interpretable outputs that support inspection of intermediate and final
    results, including annotated broadcast video, bird's-eye visualisations,
    side-by-side renderings, JSON statistics, and heatmap images.
\end{itemize}

\subsection{Non-functional Requirements}
\label{subsec:non_functional_requirements}

The non-functional requirements describe the practical and technical qualities
expected of the system beyond its core functionality.

\begin{itemize}
    \item \textbf{NFR1: Reproducibility.} The system should be executable from
    stored model weights, configuration parameters, and documented runtime
    dependencies, so that results can be reproduced consistently.

    \item \textbf{NFR2: Modularity.} The pipeline should be structured as a set
    of separable stages so that individual components, such as the tracker,
    team classifier, or pitch registration module, can be analysed, replaced,
    or improved independently.

    \item \textbf{NFR3: Robustness to broadcast conditions.} The system should
    remain usable under common broadcast challenges, including camera motion,
    partial occlusion, missed detections, and noisy pitch registration.

    \item \textbf{NFR4: Interpretability.} Intermediate and final outputs
    should be inspectable by a human observer, allowing incorrect detections,
    unstable tracks, registration failures, and tactical inference errors to be
    identified and analysed.

    \item \textbf{NFR5: Practical runtime.} The system should process complete
    match clips in a practically manageable time, even if the design prioritises
    analytical accuracy over real-time deployment.
\end{itemize}

\subsection{Success Criteria}
\label{subsec:success_criteria}

The success of the system was defined in terms of both component-level and
end-to-end behaviour.

At the component level, the detector should localise the relevant classes with
useful accuracy, the tracker should preserve identities with limited
fragmentation, the team identification stage should assign stable team labels,
and the pitch registration stage should produce sufficiently consistent
transformations for downstream pitch-space analysis. At the end-to-end level,
the pipeline should generate coherent visualisations and plausible tactical
outputs across real broadcast clips.

In practical terms, the project was considered successful if it could:

\begin{itemize}
    \item produce stable annotated video outputs showing tracked and labelled
    entities;
    \item generate bird's-eye projections that remain visually aligned with the
    broadcast scene;
    \item produce usable player heatmaps and per-player spatial statistics;
    \item estimate possession and pass events in a way that is consistent with
    manual inspection of representative clips; and
    \item identify candidate offside situations in pitch space under the
    assumptions defined later in this chapter.
\end{itemize}

These success criteria also shaped the evaluation design. Rather than judging
the project only by model-level training metrics, the system is assessed both
at the level of individual components and as an integrated pipeline, since the
ultimate objective is not only accurate reconstruction, but analytically useful
football outputs from broadcast video.

\section{System Architecture}
\label{sec:system_architecture}

The system was designed as a modular pipeline in which each stage transforms
the output of the previous stage into a representation that is more useful for
football analysis. This design was chosen because the overall task is too
complex to solve reliably as a single monolithic prediction problem: player
locations, ball position, team identity, pitch geometry, and tactical events
each require different forms of visual and geometric reasoning. A modular
architecture also supports debugging, replacement of individual components,
and component-level evaluation, which are important for reproducibility and
critical analysis.

\subsection{Pipeline Overview}
\label{subsec:pipeline_overview}

Figure~\ref{fig:pipeline_overview} shows the overall architecture of the
pipeline. The input is a broadcast football video. The output is an annotated
video together with optional bird's-eye visualisations, side-by-side videos,
heatmaps, possession and pass statistics, and detected offside candidate
events.

\begin{figure}[ht]
    \centering
    \includegraphics[width=\linewidth]{images/pipeline_overview.png}
    \caption{Overview of the proposed broadcast football video analysis
    pipeline. The system begins with object and ball detection, then applies
    tracking, team identification, pitch registration, and tactical analysis
    to produce visual and quantitative outputs.}
    \label{fig:pipeline_overview}
\end{figure}

The pipeline consists of two main phases. The first is a warm-up phase used to
fit the team identification model from player crops sampled near the beginning
of the video. The second is a per-frame processing phase in which detections,
tracks, team labels, pitch coordinates, visualisations, and statistics are
updated continuously.

\subsection{Data Flow and Intermediate Representations}
\label{subsec:data_flow_intermediate_representations}

Each frame is first processed by the main object detector, which identifies
players, goalkeepers, referees, and ball candidates. A dedicated ball detector
is then applied to improve ball localisation, since the ball is small and
frequently missed by the general detector. The resulting detections are passed
to a multi-object tracker, which assigns persistent track identifiers to
players across frames.

Team identification is performed using player crops. During the warm-up phase,
player crops are sampled from the video and used to fit an unsupervised team
classifier. During the per-frame phase, each tracked player is assigned a team
label, and temporal smoothing is applied to reduce frame-level instability.
Goalkeepers are handled separately using spatial proximity to the detected
outfield teams, while referees are assigned to a non-team class.

Pitch registration is performed by detecting pitch keypoints and estimating a
homography between the broadcast frame and a canonical pitch model. This
allows image-space detections to be transformed into pitch-space coordinates.
Pitch-space coordinates are then used to generate bird's-eye visualisations,
heatmaps, speed and distance estimates, possession statistics, pass events,
and offside candidate detections.

The main intermediate representations used by the system are:

\begin{itemize}
    \item \textbf{Bounding boxes}, representing detected players, goalkeepers,
    referees, and ball candidates in image space.

    \item \textbf{Track identifiers}, representing temporally consistent
    player identities across frames.

    \item \textbf{Team labels}, representing the inferred team membership of
    each tracked player.

    \item \textbf{Pitch keypoints}, representing detected landmarks on the
    football pitch.

    \item \textbf{Homography matrices}, representing the mapping from image
    coordinates to pitch coordinates.

    \item \textbf{Pitch-space coordinates}, representing object locations in a
    canonical coordinate system suitable for tactical analysis.
\end{itemize}

\subsection{System Outputs}
\label{subsec:system_outputs}

The system produces several forms of output so that both intermediate behaviour
and final tactical results can be inspected. The main output is an annotated
broadcast video in which detected players, goalkeepers, referees, and the ball
are visualised with track and team information. When pitch registration is
enabled, the system can also generate a bird's-eye video and a side-by-side
video combining the original broadcast view with the pitch-space
representation.

In addition to video outputs, the pipeline can export structured data. Player
and team statistics are written to JSON files, including possession estimates,
pass counts, player speed estimates, and distance travelled. Candidate offside
events are also exported as structured records, allowing them to be inspected
after processing. Player heatmaps are saved as image files and raw numerical
arrays so that the spatial distribution of each tracked player can be analysed
or reproduced later.

This separation between visual outputs and structured outputs is important
because the system must support both qualitative inspection and quantitative
evaluation. Visual outputs make it possible to diagnose errors such as missed
detections, team-label switches, or unstable registration, while structured
outputs allow the downstream tactical analysis to be evaluated more
systematically in Chapters~\ref{ch:results} and~\ref{ch:evaluation}.

\section{Datasets and Model Training}
\label{sec:datasets_model_training}

Three supervised models were trained for the system: a main object detection
model, a dedicated ball detection model, and a pitch keypoint detection model.
The main detector provides the player, goalkeeper, referee, and initial ball
detections used by the tracking pipeline. The dedicated ball detector is used
to improve ball localisation, while the pitch keypoint model provides the
landmarks required for homography estimation. Team identification does not use
a separately trained supervised model; instead, it is fitted at runtime using
player crops from the input video.

\subsection{Object Detection Dataset and Training}
\label{subsec:object_detection_dataset_training}

The main object detector was trained on the Roboflow
\texttt{football-updated-2024} dataset, version 2. This dataset contains four
classes: \textit{ball}, \textit{goalkeeper}, \textit{player}, and
\textit{referee}. The training split contained 618 images and the validation
split contained 75 images.

A YOLOv8m object detection model was fine-tuned on this dataset for 50 epochs.
The model was trained with an image size of 1280 pixels and a batch size of 4.
This model size was selected as a compromise between detection accuracy and
runtime cost, since the detector is executed throughout the full video
processing pipeline. The resulting model weights were saved as
\texttt{models/object\_detection/best.pt} and used as the main detector during
pipeline inference.

\subsection{Ball Detection Dataset and Training}
\label{subsec:ball_detection_dataset_training}

A separate ball detector was trained because the ball is substantially smaller
than the other object classes and is often affected by motion blur, occlusion,
or compression artefacts. In preliminary pipeline testing, relying only on the
ball class from the main detector produced missed detections that affected
possession, pass detection, and offside candidate generation. A dedicated model
therefore allowed the detection stage to be specialised for the most difficult
object in the scene.

The ball detector used a YOLOv8x architecture and was trained on a dedicated
Roboflow ball dataset. Training was run for 150 epochs with an image size of
1280 pixels, a batch size of 8, patience of 30 epochs, and caching enabled. The
trained weights were saved as \texttt{models/ball\_detection/best.pt}. During
inference, this detector is used alongside the main detector, with additional
temporal filtering applied later in the pipeline to reject implausible ball
jumps.

\subsection{Pitch Keypoint Dataset and Training}
\label{subsec:pitch_keypoint_dataset_training}

Pitch registration requires correspondences between points in the broadcast
image and points on a canonical football pitch. To estimate these
correspondences automatically, a YOLOv8x-pose model was trained to detect 32
canonical pitch keypoints. The model was trained using the Roboflow
\texttt{football-field-detection} dataset, version 2, which contains one
\textit{pitch} class and 32 labelled keypoints. The training split contained
255 images and the validation split contained 34 images.

The pitch keypoint model was trained for 500 epochs with an image size of 640
pixels and a batch size of 48. Mosaic augmentation was disabled because
geometric consistency between pitch markings is important for keypoint
localisation and homography estimation. The resulting model weights were saved
as \texttt{models/pitch\_keypoints/best.pt}. At inference time, predicted
keypoints are filtered by confidence before being used to estimate the
homography.

\subsection{Training Configuration}
\label{subsec:training_configuration}

All YOLO models were trained using the Ultralytics framework. The default
training configuration used automatic optimiser selection, an initial learning
rate of 0.01, final learning-rate fraction of 0.01, momentum of 0.937, weight
decay of 0.0005, warm-up for 3 epochs, automatic mixed precision, deterministic
training, and seed 0 where supported.

Table~\ref{tab:model_training_summary} summarises the supervised models used
in the system.

\begin{table}[ht]
    \centering
    \begin{tabular}{lllll}
        \toprule
        \textbf{Model} & \textbf{Architecture} & \textbf{Task} &
        \textbf{Epochs} & \textbf{Image size} \\
        \midrule
        Main detector & YOLOv8m & Player/GK/referee/ball detection &
        50 & 1280 \\
        Ball detector & YOLOv8x & Ball detection &
        150 & 1280 \\
        Pitch keypoints & YOLOv8x-pose & 32 pitch keypoints &
        500 & 640 \\
        \bottomrule
    \end{tabular}
    \caption{Summary of supervised models trained for the pipeline.}
    \label{tab:model_training_summary}
\end{table}

The trained weights are treated as external artefacts required for
reproducibility. They are stored separately from the source code and must either
be supplied with the project or regenerated using the documented training
scripts and notebooks. Model-level validation results are reported in
Chapter~\ref{ch:results}, while the effect of these models on the full pipeline
is evaluated in Chapter~\ref{ch:evaluation}.

\section{Detection and Tracking}
\label{sec:detection_tracking}

The first operational stage of the pipeline detects the visual entities required
for downstream football analysis and assigns persistent identities to players
across frames. This stage is critical because later components, including team
identification, pitch-space projection, possession estimation, heatmaps, and
offside candidate detection, all depend on the reliability of the detections
and tracks produced here.

\subsection{Main Object Detection}
\label{subsec:main_object_detection}

The main detector is a YOLOv8m object detection model trained to recognise four
classes: players, goalkeepers, referees, and the ball. During inference, each
video frame is passed through the detector at an image size of 1280 pixels. The
runtime confidence threshold is set to 0.2 and the IoU threshold is set to 0.5.
These values were chosen to retain sufficient detections for the tracker while
still removing low-confidence false positives.

The detector output for each frame consists of a set of bounding boxes, class
labels, and confidence scores. These detections are converted into
\texttt{supervision} detection objects so that they can be passed consistently
between the object detection, tracking, annotation, and tactical analysis
modules. Preserving this intermediate representation also makes it easier to
inspect and evaluate the behaviour of individual pipeline stages.

Although the main detector includes a ball class, the ball is substantially
harder to detect than players because it occupies very few pixels and is often
blurred or occluded. For this reason, the main detector is used primarily for
players, goalkeepers, and referees, while ball localisation is strengthened
using a separate dedicated detector.

\subsection{Dedicated Ball Detection and Filtering}
\label{subsec:dedicated_ball_detection_filtering}

A dedicated YOLOv8x ball detector is applied to each frame to improve the
reliability of ball localisation. This design decision was made because errors
in ball detection propagate directly into possession, pass, and offside
analysis. A missed or incorrect ball location can cause the system to assign
possession to the wrong player, miss a pass transition, or trigger an incorrect
offside candidate event.

The ball detector is run with a low initial confidence threshold
(\texttt{ball\_conf = 0.05}) so that weak ball candidates are not discarded too
early. Candidate detections are then filtered temporally. A ball candidate is
accepted only if it is sufficiently confident or spatially consistent with the
previous accepted ball position. In particular, the pipeline uses a maximum
jump threshold of 80 pixels and a minimum stronger-confidence threshold of
0.25. This prevents isolated low-confidence detections far from the previous
ball position from being selected as the ball.

When multiple ball candidates are present, the system selects the candidate
that is most consistent with the recent ball position. This filtering strategy
does not solve all cases of ball occlusion or fast motion, but it reduces the
effect of isolated false positives and produces a more stable ball trajectory
for downstream tactical reasoning.

\subsection{Multi-object Tracking}
\label{subsec:multi_object_tracking}

Player tracking is performed using tracking-by-detection. In each frame, the
main YOLO detector supplies bounding boxes to a tracker, which assigns a
persistent track identifier to each detected player. The default tracker used
in the pipeline is ByteTrack, with BoT-SORT available as an alternative
configuration. The tracker is selected through the runtime tracker
configuration, allowing the same pipeline to be evaluated with different
tracking algorithms.

ByteTrack was selected as the default tracker because it is designed to use
both high-confidence and lower-confidence detections during association. This
is useful in broadcast football footage, where players are frequently
partially occluded and their detection confidence may temporarily fall even
when they remain visible. By retaining lower-confidence detections for a second
association step, ByteTrack can reduce unnecessary track fragmentation in
crowded scenes.

The tracker returns detections with associated track identifiers. These track
identifiers are preserved in the pipeline and are used by later stages to
accumulate player positions, smooth team labels, generate heatmaps, and compute
player-level statistics. Detections without valid track identifiers are not
used for player-level temporal statistics, since they cannot be linked
reliably across frames.

\subsection{Track Handling and Failure Cases}
\label{subsec:track_handling_failure_cases}

The tracking stage is vulnerable to several common failure cases. Occlusion can
cause a player to disappear temporarily, leading to a broken track. Similar
player appearance can cause identity switches, particularly when players cross
or cluster around the ball. Camera motion can also affect the apparent movement
of detections and make association more difficult.

The pipeline mitigates these issues in two ways. First, the tracker itself
uses temporal association to preserve identities across short interruptions.
Second, downstream stages avoid relying on a single frame wherever possible.
For example, team labels are smoothed across a fixed track history, and
movement statistics include a maximum-speed gate to reduce the effect of
spurious projected movements.

Nevertheless, tracking errors cannot be fully eliminated at this stage. A
fragmented or incorrect track can affect player heatmaps, distance estimates,
team-level possession, and offside candidate detection. For this reason, the
evaluation considers both the component-level behaviour of the tracker and the
effect of tracking errors on the integrated pipeline.

\section{Team Identification}
\label{sec:method_team_identification}

After player detections have been assigned track identifiers, the next stage is
to infer the team membership of each tracked player. This is required for
team-level outputs such as possession statistics, pass events, offside
candidate detection, and team-coloured visualisations. The system performs
team identification without match-specific manual labels by fitting an
unsupervised classifier from player crops sampled from the input video.

\subsection{Warm-up Sampling and Player Cropping}
\label{subsec:warmup_sampling_player_cropping}

Team identification begins with a warm-up phase at the start of each video.
During this phase, frames are sampled from the first 20 seconds of footage using
a stride of 30 frames. The main detector is applied to these sampled frames, and
only detections classified as players are retained. Very small crops are
discarded using a minimum bounding-box size of 35 by 35 pixels, since such
detections usually contain too little visual information for reliable team
assignment.

For each retained player detection, the system extracts a torso crop rather
than using the full bounding box. The crop is taken from approximately 15\% to
65\% of the bounding-box height. This focuses the representation on the shirt
region, which usually contains the strongest team-specific visual signal, while
reducing noise from shorts, legs, grass, and nearby players. A maximum of 800
warm-up crops is used, and at least 50 crops are required before the classifier
is fitted.

\subsection{Embedding and Clustering}
\label{subsec:embedding_clustering}

The retained torso crops are used to fit a team classifier. The classifier
extracts visual embeddings from each crop using a SigLIP-based image encoder,
then reduces the embedding dimensionality using UMAP and clusters the reduced
representations using K-means with \(k=2\). The two clusters correspond to the
two outfield teams.

This approach was selected because it avoids requiring labelled team examples
for every new match. A supervised team classifier would need new annotations
whenever kits, lighting conditions, competitions, or broadcast styles changed.
By fitting the classifier from the input video itself, the system can adapt to
the visual appearance of the teams in the current match.

During the per-frame pipeline, each tracked player crop is passed through the
same embedding and clustering pipeline to obtain a predicted team label. The
output is a binary team assignment for outfield players. These assignments are
then stabilised using temporal smoothing before being used by downstream
analytics.

\subsection{Temporal Smoothing and Role Handling}
\label{subsec:temporal_smoothing_role_handling}

Frame-level team predictions can fluctuate because of occlusion, motion blur,
partial crops, lighting variation, or players turning away from the camera. To
reduce this instability, the system applies track-level majority-vote smoothing.
For each track identifier, the most recent team predictions are stored in a
fixed-length history of 30 predictions. The smoothed team label is then taken as
the majority label within this history.

This design reduces the effect of isolated incorrect predictions while still
allowing the team label to update if a track was initially misclassified.
Smoothing is especially important because team labels are used in later
components where a single-frame error can affect possession, pass, or offside
logic.

Goalkeepers and referees require separate handling. Referees are assigned to a
non-team class because they do not contribute to team possession or offside
analysis. Goalkeepers are detected as a separate object class and are assigned
to the team whose outfield-player centroid is nearest in image space. This
spatial rule is used because goalkeeper kits often differ from the outfield
players' kits, making appearance-based clustering less reliable for this role.

The resulting representation assigns each detection either to team 0, team 1,
or a non-team class. This representation is used consistently by the
visualisation, possession, pass, and offside modules.

\section{Pitch Registration}
\label{sec:method_pitch_registration}

Pitch registration maps image-space detections from the broadcast frame onto a
canonical football pitch. This stage is required because image coordinates are
affected by camera position, zoom, and perspective distortion. Tactical outputs
such as heatmaps, speed, distance, possession proximity, and offside candidate
detection therefore need to be computed in pitch space rather than directly in
the broadcast image.

\subsection{Pitch Keypoint Detection}
\label{subsec:pitch_keypoint_detection}

Pitch landmarks are detected using a YOLOv8x-pose model trained to localise 32
canonical pitch keypoints. For each processed frame, the model predicts a set
of keypoint coordinates and confidence scores. Keypoints with confidence below
0.5 are discarded before homography estimation. This filtering step reduces the
effect of uncertain keypoint predictions, which can otherwise produce unstable
or incorrect pitch transformations.

A homography is estimated only when at least six valid keypoints are available.
Although four point correspondences are theoretically sufficient to estimate a
homography, requiring more than four points makes the transformation less
dependent on any single detected landmark. Frames with too few reliable
keypoints reuse the most recent valid transformation where possible, rather
than forcing an unstable new estimate.

\subsection{Homography Estimation and Coordinate Transformation}
\label{subsec:homography_coordinate_transformation}

For each valid set of pitch keypoints, the system estimates a homography
between the detected image-space landmarks and their corresponding locations on
a canonical pitch model. The homography provides a projective transformation
from image coordinates to pitch coordinates. Once this transformation is
available, detected object positions can be expressed in pitch space.

For players, goalkeepers, and referees, the image-space position is taken as
the bottom-centre point of the bounding box, since this approximates the point
where the player contacts the ground. For the ball, the centre of the bounding
box is used. These points are then transformed into pitch coordinates using the
estimated homography.

The pitch coordinate system is represented in centimetres, with a pitch length
of 12,000 cm. Using metric-like pitch units allows later stages to express
distance thresholds and spatial rules in interpretable terms. For example, the
possession and offside modules use distances in centimetres rather than pixel
distances, making the thresholds less dependent on camera zoom.

\subsection{Caching and Robustness Measures}
\label{subsec:caching_robustness_measures}

Estimating a new pitch transformation on every frame is computationally
expensive and can also introduce visual jitter when keypoint predictions vary
slightly between consecutive frames. To reduce this problem, the pipeline
recomputes the pitch transformation every 15 frames and caches the most recent
valid transformation between updates.

This stride-cached design is a compromise between accuracy, stability, and
runtime. Updating the homography too rarely can make the transformation lag
behind camera motion, while updating it every frame can amplify keypoint noise.
A stride of 15 frames provides regular updates while reducing unnecessary
recomputation.

Pitch registration remains one of the most important sources of downstream
error. If the homography is inaccurate, then all pitch-space outputs derived
from it may be affected, including heatmaps, speed and distance estimates,
possession proximity, and offside geometry. For this reason, registration
quality is inspected both visually through bird's-eye outputs and quantitatively
where keypoint validation metrics are available.

\section{Tactical Output Generation}
\label{sec:tactical_output_generation}

Once detections have been tracked, assigned team labels, and transformed into
pitch coordinates, the system generates tactical outputs from the resulting
structured representation. These outputs are intended to make the broadcast
video more analytically useful by converting low-level detections into visual
summaries, player statistics, team possession estimates, pass events, and
candidate offside situations.

\subsection{Bird's-eye Visualisation}
\label{subsec:birdeye_visualisation}

The bird's-eye visualisation represents the detected players, goalkeepers,
referees, and ball on a canonical top-down pitch. For each frame with a valid
pitch transformation, the image-space object positions are projected into pitch
coordinates and rendered on the pitch model. Players are coloured according to
their inferred team labels, referees are shown as a separate non-team class,
and the ball is rendered using its filtered pitch-space position when
available.

This visualisation provides an interpretable way to inspect whether the
detection, team assignment, and pitch registration stages are behaving
coherently. It is also used to generate side-by-side videos, where the original
broadcast frame and the corresponding pitch-space representation are displayed
together. This makes it easier to identify errors such as incorrect team
labels, unstable homographies, or ball locations that do not match the
broadcast footage.

\subsection{Heatmaps, Speed and Distance}
\label{subsec:heatmaps_speed_distance}

Player heatmaps are generated by accumulating each tracked player's pitch-space
positions over time. The pitch is divided into a fixed grid of
\(480 \times 280\) cells, and each valid player position contributes to the
corresponding grid location. Heatmaps are generated only when at least 40 valid
samples are available for a track, reducing the likelihood of producing
misleading visualisations from short or fragmented tracks.

To improve readability, the accumulated heatmap is smoothed using a Gaussian
blur with kernel size 31 and a gamma correction of 0.22. The final heatmap is
saved both as an image and as a raw numerical array, so that the visual output
can be inspected while the underlying positional distribution remains
available for further analysis.

Speed and distance estimates are computed from changes in player pitch-space
position over time. Consecutive valid positions for each track are used to
estimate displacement, which is accumulated to produce total distance
travelled. Speed is derived from displacement and frame timing. Since small
errors in pitch registration can produce unrealistically large apparent jumps,
the system applies a maximum-speed gate of 40 km/h. Movements above this
threshold are treated as noise and excluded from the accumulated statistics.

\subsection{Possession and Pass Detection}
\label{subsec:possession_pass_detection}

Possession is estimated using the pitch-space distance between the filtered
ball position and nearby players. For each frame, the system identifies the
player closest to the ball. A player is considered a possession candidate only
if the distance to the ball is below 350 cm. To reduce noise from brief
misdetections or transient ball movements, possession must be stable for at
least three consecutive frames before being confirmed.

Team possession is then obtained from the confirmed player's team label. If no
player satisfies the proximity condition, the frame is treated as having no new
confirmed possession event. This conservative design reduces the risk of
assigning possession during periods when the ball is airborne, occluded, or
incorrectly detected.

Pass events are detected from stable changes in possession. A pass is recorded
when confirmed possession changes from one player to another in a temporally
stable way. This rule-based definition is intentionally simple: it is designed
to identify likely possession transfers rather than to provide a complete
semantic classification of football actions. As a result, the pass output is
best interpreted as a set of candidate pass events that require evaluation
against manually inspected clips.

\subsection{Offside Candidate Detection}
\label{subsec:offside_candidate_detection}

Offside candidate detection is performed in pitch space using the detected
players, inferred team labels, ball position, and estimated attacking
direction. The attacking direction is inferred automatically from accumulated
team \(x\)-centroids once at least 30 valid team-position samples are
available, rather than being manually specified for each video.

For each frame with sufficient information, the system estimates the offside
line from the second-last defender's pitch-space position. If only one defender
is available, a fallback line is computed using the available defensive player,
although this case is less reliable. An attacker is treated as a candidate
offside player if they are beyond the defensive line by more than a 30 cm
margin at the relevant moment and are spatially close enough to the ball event
context to be considered involved in play.

The output of this module is not intended to replicate full official offside
judgement. The official law depends on contextual factors such as the exact
moment the ball is played, active involvement, deliberate play by defenders,
and body-part eligibility. The system instead detects candidate offside
situations under simplified geometric assumptions. Candidate events are saved
as structured JSON records and visualised using an overlay for a fixed display
duration of 50 frames, allowing them to be inspected during evaluation.

\section{Evaluation Methodology}
\label{sec:evaluation_methodology}

The evaluation was designed to assess both the individual components of the
pipeline and the behaviour of the integrated system. This distinction is
important because strong component-level performance does not necessarily imply
that the final tactical outputs are reliable. For example, a detector may
achieve high validation accuracy while still missing the ball in key moments,
or a pitch registration model may perform well on isolated frames while
producing unstable transformations during camera motion. The evaluation
therefore considers not only model-level metrics, but also the extent to which
the complete pipeline produces coherent and useful football analysis outputs.

\subsection{Component-level Evaluation}
\label{subsec:component_level_evaluation}

The first level of evaluation assesses the supervised models used by the
pipeline. The main object detector, dedicated ball detector, and pitch keypoint
detector are evaluated using their validation or test splits, with metrics
appropriate to each task.

For the object detection models, performance is measured using precision,
recall, mAP@50, and mAP@50--95. These metrics assess both whether the relevant
objects are detected and how accurately their bounding boxes are localised. The
dedicated ball detector is evaluated separately because ball localisation is a
known bottleneck for downstream possession, pass, and offside analysis.

For the pitch keypoint model, evaluation considers both bounding-box and
keypoint performance where available. The keypoint metrics are particularly
important because pitch registration depends on the geometric accuracy of the
detected landmarks rather than only on whether the pitch itself is detected.

Tracking is evaluated by comparing the stability and continuity of player
tracks across representative clips. Where manual annotations are available,
track fragmentation, identity switches, and track continuity are assessed
directly. Where dense tracking ground truth is not available, the tracker is
evaluated through visual inspection, track-count consistency, and its effect on
downstream outputs such as heatmaps and distance estimates.

\subsection{End-to-end Pipeline Evaluation}
\label{subsec:end_to_end_pipeline_evaluation}

The second level of evaluation assesses the complete pipeline on broadcast
football clips. This evaluation focuses on whether the system produces coherent
outputs after all components have been combined. The main outputs considered
are annotated broadcast videos, bird's-eye visualisations, side-by-side videos,
heatmaps, statistics files, and offside event records.

End-to-end evaluation is performed using representative clips containing common
broadcast conditions, including camera panning, zooming, player occlusion, ball
movement, and changes in pitch visibility. The outputs are inspected to assess
whether detected entities remain visually consistent, whether team labels are
stable, whether pitch-space projections remain aligned with the broadcast view,
and whether the generated statistics are plausible.

Where manual labels are produced, tactical outputs are evaluated against
manually inspected ground truth. Possession estimates are compared with the
observed player or team in control of the ball, pass events are compared with
manually identified possession transfers, and offside candidates are compared
with visible attacking situations in the clip. This provides a more meaningful
assessment of tactical usefulness than model-level metrics alone.

\subsection{Ablation Studies}
\label{subsec:ablation_studies}

Ablation studies are used to evaluate the contribution of key design decisions.
Each ablation changes one component or configuration while keeping the rest of
the pipeline fixed, allowing the effect of that design choice to be isolated.

The planned ablations are:

\begin{itemize}
    \item \textbf{Dedicated ball detector.} The pipeline is compared with and
    without the dedicated ball detector to determine whether the specialised
    model improves possession, pass, and offside outputs.

    \item \textbf{Tracker choice.} ByteTrack and BoT-SORT are compared on the
    same clips to assess differences in track continuity, identity stability,
    and downstream statistics.

    \item \textbf{Team-label smoothing.} Raw frame-level team predictions are
    compared with track-level majority-vote smoothing to assess whether
    smoothing reduces label instability.

    \item \textbf{Pitch registration stride.} Different homography update
    intervals are compared to measure the trade-off between runtime,
    projection stability, and responsiveness to camera motion.

    \item \textbf{Speed filtering.} Distance and speed estimates are compared
    with and without the maximum-speed gate to determine how much unrealistic
    movement is introduced by registration noise.
\end{itemize}

These ablations are included because the project contains several design
choices that could otherwise appear arbitrary. Comparing alternatives makes it
possible to justify the final configuration and identify which components have
the greatest effect on end-to-end reliability.

\subsection{Runtime Evaluation}
\label{subsec:runtime_evaluation}

Runtime evaluation measures the practical cost of processing broadcast video.
The system records total processing time and, where possible, the time spent in
major pipeline stages such as detection, ball detection, pitch registration,
team assignment, rendering, and tactical output generation.

Runtime is reported as processing time per frame and effective frames per
second. Additional comparisons are made for configurations that change runtime
substantially, such as different pitch registration strides or optional
bird's-eye and side-by-side rendering. This allows the evaluation to distinguish
between methods that improve output quality and methods that are too expensive
for practical use.

The aim is not necessarily to achieve real-time performance, since the project
prioritises analytical quality and reproducibility. However, runtime remains an
important non-functional requirement because a system that is too slow to run on
complete clips would have limited practical value.

\subsection{Error Analysis}
\label{subsec:error_analysis}

Error analysis is used to examine where the pipeline fails and how errors
propagate between stages. This is necessary because the final tactical outputs
depend on several earlier components. A missed ball detection can affect
possession and pass detection; an identity switch can distort heatmaps and
distance estimates; an incorrect team label can affect possession and offside
logic; and an unstable homography can corrupt all pitch-space measurements.

The evaluation therefore records representative failure cases and groups them
by cause. The main categories considered are missed or false detections, ball
tracking errors, track fragmentation, team-label switches, pitch registration
failures, and tactical inference errors. For each category, the analysis
identifies the likely cause, the downstream effect, and any mitigation used in
the pipeline.

This error-focused evaluation is important because the project is an integrated
system rather than a single model. The objective is not only to report whether
the outputs are correct, but also to explain why they succeed or fail under
different broadcast conditions. This provides the basis for the critical
evaluation in Chapter~\ref{ch:evaluation}.

\section{Implementation and Reproducibility}
\label{sec:implementation_reproducibility}

The pipeline was implemented in Python, using a modular code structure so that
individual stages could be configured, tested, and replaced independently. The
main entry point is \texttt{src/track\_video\_supervision.py}, which parses the
runtime configuration and calls the pipeline orchestration code in
\texttt{src/pipeline/run.py}. This organisation separates high-level execution
from the component-specific modules used for detection, team identification,
pitch registration, visualisation, statistics, and offside detection.

\subsection{Software Environment}
\label{subsec:software_environment}

The implementation uses Python 3.11. The main libraries are Ultralytics for
YOLO model training and inference, Supervision for detection handling and video
annotation utilities, OpenCV for image processing and video writing, NumPy for
numerical operations, scikit-learn and UMAP for clustering-related operations,
and Hugging Face Transformers for the SigLIP-based team embedding model. The
Roboflow inference tools are used for the pitch model integration, and the
Roboflow sports library is used for the team classification pipeline.

The project dependencies are recorded in \texttt{requirements.txt}. In addition,
the Roboflow sports library is installed as an editable dependency from its
source repository. The environment also requires a Roboflow API key and field
model identifier, which are specified through environment variables. An example
environment file is provided to document the required variables without exposing
private credentials.

\subsection{Configuration and Model Weights}
\label{subsec:configuration_model_weights}

The pipeline is controlled through runtime parameters rather than hard-coded
settings. Important configuration values include the detector image size,
detection confidence threshold, IoU threshold, tracker type, pitch registration
stride, keypoint confidence threshold, ball filtering thresholds, team
smoothing window, heatmap settings, and offside detection parameters.

The default inference configuration uses an image size of 1280 pixels, detector
confidence threshold of 0.2, IoU threshold of 0.5, and ByteTrack as the tracker.
The team classifier samples the first 20 seconds of video with a stride of 30
frames and applies a 30-frame smoothing window to team predictions. Pitch
registration is updated every 15 frames, using keypoints with confidence at
least 0.5 and requiring at least six valid keypoints before estimating a
homography.

The trained model weights are stored separately from the source code. The
required weights are the main object detector
\texttt{models/object\_detection/best.pt}, the dedicated ball detector
\texttt{models/ball\_detection/best.pt}, and the pitch keypoint detector
\texttt{models/pitch\_keypoints/best.pt}. To reproduce the full system, these
weights must either be supplied with the project artefacts or regenerated using
the documented training scripts and notebooks.

\subsection{Running the Pipeline}
\label{subsec:running_pipeline}

The pipeline can be executed from the command line using the provided run
scripts. The standard execution path processes an input video, writes an
annotated output video, and optionally generates bird's-eye videos,
side-by-side videos, heatmaps, statistics files, pitch debug videos, and
offside event records. The main generated outputs are written to the
\texttt{outputs/} directory.

For each input video, the core output is an annotated video showing detections,
track identifiers, team labels, the ball marker, and offside overlays where
enabled. Optional outputs include a bird's-eye video, a side-by-side
broadcast-and-pitch view, JSON statistics, JSON offside-event records, and
per-player heatmap files. The statistics output records team possession,
passes, player speed estimates, and distance travelled, while the heatmap
directory stores both image files and raw numerical arrays.

This output structure was chosen to support both visual inspection and later
evaluation. Videos make it possible to inspect whether the integrated pipeline
is behaving plausibly, while JSON and numerical outputs allow quantitative
comparison with manual labels and ablation results.

\subsection{Reproducibility Limitations}
\label{subsec:reproducibility_limitations}

Several limitations affect exact reproducibility. First, the trained model
weights are external artefacts and are not stored directly inside the source
code repository. Reproducing the system therefore requires either access to the
trained weights or re-training the models from the documented datasets and
training configuration.

Second, some datasets and hosted models are accessed through Roboflow services,
so reproduction depends on access to the relevant workspaces, dataset versions,
and API credentials. Third, although YOLO training used deterministic settings
and a fixed seed where supported, the runtime team classifier is fitted from
sampled video crops and may vary if the underlying embedding, UMAP, or
clustering behaviour is not fully deterministic.

Finally, the current implementation is primarily an experimental research
pipeline rather than a fully packaged software product. It provides scripts,
configuration parameters, model artefact locations, and documented outputs, but
does not include a complete automated unit-test suite. This limitation is
addressed in the evaluation through component-level validation, visual
inspection, ablation studies, and end-to-end testing on representative clips.

\section{Chapter Summary}
\label{sec:methodology_summary}

This chapter described the methodology used to construct the proposed
broadcast football video analysis pipeline. The system was designed around a
set of functional and non-functional requirements, with the central objective
of producing tactical outputs from a single broadcast video without relying on
commercial tracking infrastructure.

The methodology used a modular architecture in which object detection, ball
detection, multi-object tracking, team identification, pitch registration, and
tactical output generation are implemented as separable stages. This structure
allows each component to be configured and evaluated independently, while still
supporting end-to-end analysis of the complete pipeline.

The chapter also described the datasets, model training procedures, runtime
configuration, and reproducibility measures used by the system. Particular
attention was given to design decisions that affect downstream reliability,
including the use of a dedicated ball detector, track-level team-label
smoothing, stride-cached pitch registration, and pitch-space rules for
possession, pass, and offside candidate detection.

The following chapter reports the results produced by this methodology,
including model-level validation results, pipeline outputs, tactical
statistics, visualisations, and runtime measurements. Chapter~\ref{ch:evaluation}
then evaluates these results against the requirements and success criteria
defined in this chapter.
